% TOP_PROBLEM: {"id": 3100019, "title": "The \"cycle double cover conjecture\" states that every bridgeless graph contains a set of cycles which cover each edge of", "category_id": 2, "category": {"id": 2, "name": "combinatorics", "display_name": "Combinatorics", "description": "Counting problems, graph theory, discrete structures.", "slug": "combinatorics", "order_index": 2, "created_at": "2026-07-31T15:26:25.671Z"}, "status": "partially_solved", "problem_number": "AMR-030-0019", "set_id": 13}
% FIRST_POSTED: 2026-10-02
% ENTRY_KIND: statement_only
% UnsolvedMath ID 3100019; source catalogue code AMR-030-0019
% !TeX program = lualatex
% Definition and sources only; this file is not an LLM solution attempt.
\documentclass[11pt,a4paper]{article}
\usepackage[margin=25mm]{geometry}
\usepackage{fontspec}
\setmainfont{lmroman10-regular.otf}[BoldFont=lmroman10-bold.otf,
 ItalicFont=lmroman10-italic.otf,BoldItalicFont=lmroman10-bolditalic.otf]
\usepackage{amsmath,amssymb,mathtools}
\usepackage{unicode-math}
\setmathfont{latinmodern-math.otf}
\AtBeginDocument{\let\setminus\smallsetminus}
\usepackage{xurl,hyperref}
\hypersetup{colorlinks=true,urlcolor=blue}
\setcounter{secnumdepth}{0}
\setlength{\parindent}{0pt}
\setlength{\parskip}{5pt}
\setlength{\emergencystretch}{3em}
\begin{document}
\section*{The "cycle double cover conjecture" states that every bridgeless graph contains a set of cycles which cover each edge of}\label{3100019}
{\small UnsolvedMath ID 3100019 · AMR-030-0019 · Combinatorics · AMR Open Problem Lists}
\subsection{Definitions and mathematical statement}
Let $G=(V,E)$ be a finite undirected bridgeless graph: deleting any
one edge does not increase the number of connected components.
A cycle is a connected $2$-regular subgraph. For a simple graph
it has length at least three. The equivalent loopless multigraph
formulation permits a two-edge cycle formed by parallel edges.

A cycle double cover is a finite multiset $\mathcal C$ of cycles
such that
\[
            \sum_{C\in\mathcal C}\mathbf 1_{\{e\in E(C)\}}=2
                       \qquad\text{for every }e\in E.
\]
Cycle occurrences are counted with multiplicity. The cycle double
cover conjecture asserts that every finite bridgeless graph has
such a cover.

Repeated cycles are allowed: for example, a graph consisting of
one cycle is covered by taking that cycle twice. The cycles may
overlap in vertices and edges, but the total use of each edge must
be exactly two. Arbitrary fractional weights on cycles would give
a different relaxation. A bridge cannot lie on any cycle, so the
bridgeless hypothesis is necessary. Disconnected graphs are
included, and their components can be treated independently;
isolated vertices impose no covering requirement. The problem
asks for existence for all graphs, without assuming planarity,
regularity, or a particular edge coloring.
\subsection{Short English statement}
Does every finite bridgeless graph admit a multiset of cycles covering every edge exactly twice?
\subsection{Sources}
\begin{itemize}
\item[S1] ULAM.AI and the UnsolvedMath Contributors, supplied problems.json, record 3100019 (AMR-030-0019), AMR Open Problem Lists. Catalogue identity and source transcription; CC BY 4.0.
\url{https://huggingface.co/datasets/ulamai/UnsolvedMath}.
\item[S2] Cooper - Combinatorial Problems I Like (2020 snapshot), Graphs, Hypergraphs, Set Systems, Designs, source-order bullet 19. Original source indexed by AMR-030-0019.
\url{https://people.math.sc.edu/cooper/combprob.html}.
\end{itemize}
\end{document}
